\documentclass[10pt,letterpaper]{article}
\usepackage{cornell-notes}

\title{Clause 29.11.06: Class Template Zoned Traits}
\author{}
\date{}
\setDocTitle{Clause 29.11.06: Class Template Zoned Traits}
\setDocSubtitle{C++ 2024}
\setDocOwner{C++ 2024 Cornell Notes}
\setDocDate{\today}
\setCornellCollection{C++ 2024 Cornell Notes}
\setCornellUnitType{Clause}
\setCornellUnitNumber{29.11.06}
\setCornellUnitTitle{Class Template Zoned Traits}
\hypersetup{pdftitle={Clause 29.11.06: Class Template Zoned Traits},pdfsubject={Cornell notes},pdfauthor={}}

\begin{document}
\maketitle
\begin{CornellOverviewBox}{Overview}
\textbf{Classification:} Standard library - Normative\\[2pt]
\textbf{Learning objective:} Understand the rules, terminology, and semantic consequences associated with Class template zoned\_traits.
\end{CornellOverviewBox}\begin{CornellSummaryBox}{Summary}
{\sffamily\bfseries\color{Accent} Summary}\par\vspace{3pt}Section 29.11.6 focuses on Class template zoned\_traits. Apply the specified interface together with its mandates, effects, result conditions, exception behavior, and complexity constraints. When applying it, preserve the distinction between mandatory requirements, recommendations, permissions, and behavior deliberately left undefined, unspecified, or implementation-defined.
\end{CornellSummaryBox}

\end{document}
